\section{实验目的与任务}
CIFAR-10包含飞机、汽车、鸟等10类低分辨率自然图像\cite{krizhevskyCifar}。本实验围绕课程A2要求，构建可训练的卷积神经网络，研究Dropout与多种归一化对泛化的影响，通过随机搜索和分层交叉验证选择超参数，并对性能提升与网络局限开展验证。

卷积通过局部连接与参数共享利用图像的空间结构\cite{lecun1998gradient}。与将图像展平的多层感知机相比，它提供了更合适的结构先验。本报告据此区分三个问题：基础CNN能否稳定学习；正则化与归一化在固定预算下如何改变验证表现；在提高分类准确率之外，能否改善深度退化、校准和平移稳定性。

实验共定义107次正式训练，包含多随机种子对照、24组随机搜索、四候选的五折验证及最终重训。作者此前A1报告的NumPy MLP测试准确率50.39\%作为历史参考。
